% !TeX root = ../nc-cubic-non-as.tex
\subsection{既約種数一曲線型の内在的認識}
\label{subsec:ou-elliptic-recognition}

六点爆発の一般位置を仮定せず、反標準spherical functorから
超曲面表示を回復する。滑らかなadmissible elliptic decuple
\cite[Definition~7.1, Lemma~7.4]{AOU2024}の型に加えて、
節点・尖点三次曲線による表示にも適用する。
任意のOU三次型がこのような表示を持つとは仮定しない。

\begin{theorem}[既約種数一曲線のhelixを持つ内在的三次候補]
\label{thm:ou-integral-elliptic-recognition}
$E$を標数零の代数閉体上の既約被約射影算術種数一曲線
（滑らかな楕円曲線、節点三次曲線または尖点三次曲線）とし、
$(L_1,\ldots,L_9)$を次数$0,1,2$の各三本の線束からなり、
各ブロック内で相異なるものとする。
$\Lambda$をそのdirected endomorphism algebraとし、
$\mathcal C=\qgr A$を対応するOU三次型のnoetherian helix商圏とする。
標準の比較$D^b(\mathcal C)\simeq\Perf\Lambda$が成立し、
標準対象がfull very strong helix $(P_i)$に移り
$P_{i+9}\simeq sP_i$と仮定する
（滑らかな場合はAOUのadmissible elliptic decupleの設定）。
$(\mathcal C,\mathcal O)$が定義\ref{def:intrinsic-dp}を満たすなら、
与えられた標識の反標準環$R$について
\[
 R=Q(R)/(g),\qquad H_{Q(R)}(t)=(1-t)^{-4},
\]
ここで$Q(R)$は標準二次Koszul AS四次元正則環、
$g$は正則正規三次元である。さらに$R$には正則中心一次元$z$があり、
$R/zR$は$E$上の次数三線束と$\Pic^0(E)$に自明に作用する
自己同型によるtwisted homogeneous coordinate ringである。
\end{theorem}

証明では、特異曲線の場合に$\Perf E$と$D^b(\coh E)$を区別する。
$E$はGorensteinかつ$\omega_E\simeq\mathcal O_E$であり、
$\Perf E$のSerre関手は$[1]$である。
次数差一の二線束の非零射の余核は滑らかな一点の構造層：
実際、次数一線束の非零切断の零因子は長さ一の有効Cartier因子で、
特異点では極大イデアルが主イデアルにならない。
この点と一方の線束から、点による正負の捩りの完全列を用いて
任意の整数冪の捩りを得る。これらは$\Perf E$を厚く生成する。
例えば任意のベクトル束に対する十分負の捩りによる二段の全射分解は
$\operatorname{Ext}^2$の消滅（第一変数がベクトル束）により
もとの束をその二項複体の直和因子として与える。

曲線へのfunctorを$F:\Perf\Lambda\to\Perf E$とし、
$F(P_i)=L_i$とする。これはdirected dg categoryの埋込みからの
導来拡大である。右随伴$J$は$\Perf E$上の制限functorであるが、
\emph{全ての}$D^b(\coh E)$に
\[
 J(M)=\operatorname{RHom}_E\bigl(\bigoplus_{i=1}^9 L_i,M\bigr)
 \quad\text{（directed algebraへの制限）}
\]
と定義して拡張できる。各$L_i$は完全なのでこの複体は有界有限次元、
$\Lambda$は有限大域次元なのでその値は$\Perf\Lambda$に属する。
$F\dashv J$は生成射影上の随伴から全ての$\Perf\Lambda$について従う。
特異曲線の$D^b(\coh E)$自体のSerre関手は仮定していない。

$\Perf E$上でOU\cite[(3)--(5), p.~2]{OUSpherical2020}の
spherical collectionの構成を適用すると、
\[
 s^{-1}\xrightarrow{\epsilon}\mathrm{id}\longrightarrow JF,\qquad
 D:=\operatorname{Cone}(FJ\longrightarrow\mathrm{id})
   =T_{L_1}\cdots T_{L_9}.
\]
ここで$T_L$はspherical twistである。
右随伴の右随伴について$G\simeq Fs[-1]$（$\Perf$上のSerre双対性）
および$G\simeq D^{-1}F[1]$
（\cite[Definition~5.2, (5.12)]{AnnoLogvinenko2017}）から、
$h=D^{-1}[2]$と置けば
\begin{equation}\label{eq:ou-restriction-commute}
 Fs\simeq hF.
\end{equation}
各twistは完全なspherical対象$L_i$に沿うものであり、
特異曲線でも$D^b(\coh E)$上の自己同値を与える
（\cite[Sections~3--4]{BurbanKreu2006}）。
従って$h$もこの大きい導来圏へ延長する。
以下では上の同型を反復してorbit積と両立する同定を選ぶ。

\begin{lemma}\label{lem:ou-elliptic-shift}
$h$は$\coh E$の完全自己同値であり、
$h(X)\simeq \alpha^*X\otimes N$と書ける。
$N$は次数三の線束で、$\alpha^*|_{\Pic^0(E)}=\mathrm{id}$。
滑らかな場合$\alpha$は平行移動である。
\end{lemma}
\begin{proof}
$(r,d)$に対する次数$e$の線束のtwistの作用は
\[
 T_e=\begin{pmatrix}1+e&-1\\e^2&1-e\end{pmatrix}.
\]
従って$h^{-1}=D[-2]$は
$T_0^3T_1^3T_2^3=\left(\begin{smallmatrix}1&0\\-3&1\end{smallmatrix}\right)$
を誘導する。Burban--Kreu{\ss}ler
\cite[Corollary~5.8]{BurbanKreu2006}により、
次数$-3$のtensor積を取り除いた$h^{-1}$は、
曲線自己同型・次数零線束・偶数シフトの合成となる。
これは既約特異種数一曲線にも適用される。
偶数シフトを次の実計算で決定する。

最後のブロックの線束$L=L_9$について、同じブロックのtwistは
$L$を保つ。次のブロックの三本（次数一）を$M_i$とすれば
評価$\bigoplus_i M_i\to L$は全射である：各成分の零点は
滑らかな相異なる一点である。核$K$はrank二、次数一のベクトル束。
$K$は安定である。実際、次数一以上のrank一torsion-free部分層$H$
からいずれかの$M_i$への非零射は単射であり、次数比較から
$H\simeq M_i$でなければならない。相異なる他の成分への射は零なので、
この一成分は評価の核には入らない。従って
$T_{M_1}T_{M_2}T_{M_3}L=K[1]$。
最初のブロックの三本$U_i$について安定性とSerre双対性
（ベクトル束に対する双対性）から
$\dim\operatorname{Hom}(U_i,K)=1$、
$\operatorname{Ext}^1(U_i,K)=0$である。
評価の任意の二成分は一般点で独立である：
同じ飽和rank一部分層に入れば、その次数は安定性から零以下であり、
二本の$U_i$からの単射と次数比較により両方と同型となって矛盾する。
二成分の行列式は次数一の線束の非零切断であり、
一成分を共通とする二つの行列式の零点は滑らかかつ相異なる。
従って$\bigoplus_iU_i\to K$も全射、その核は次数$-1$の線束となる。
同じブロックの線束は完全直交するため、
twistの合成を上の同時評価で計算できる。
よって$h^{-1}L$はシフトのない線束である。
上記の偶数シフトは零となり、$h=\alpha^*(-)\otimes N$、
$\deg N=3$を得る。

自己同型の$\Pic^0$上の作用を確認するには
$K_0(\Perf E)$を用いる。特異曲線の$G_0(E)$で代用してはならない。
任意の次数零線束$Q$について
$x=[Q]-[\mathcal O_E]$のrank・次数は零なので、
全ての$T_{L_i}$は$x$を固定する。
$h$の幾何的表示の下でdeterminantを取ると
$\alpha^*Q\simeq Q$である。これより主張を得る。
\end{proof}

\begin{lemma}[曲線のheartと次元]\label{lem:ou-elliptic-heart}
$J$は標準heartに関して完全で、$\coh E$上でconservativeである。
任意の$M\in\coh E$の像$JM$は$\mathcal C$の
Gabriel--Krull次元一以下の対象である。
\end{lemma}
\begin{proof}
$P=\bigoplus_{i=1}^9P_i$、$L_\bullet=FP$と置き、
周期をまとめたorbit環を
$B=\bigoplus_{n\geq0}\operatorname{Hom}(P,s^nP)$とする。
これは九つの添字をまとめた通常のgraded helix algebraで、
graded module圏とtorsion商は周期を展開したものと同一視する。
まず曲線側の
\[
 B_E=\bigoplus_{n\geq0}\operatorname{Hom}_E(L_\bullet,h^nL_\bullet)
\]
と比較する。三角と(\ref{eq:ou-restriction-commute})により、
$n\geq1$では
\[
 0\to B_{n-1}\xrightarrow{\epsilon} B_n\to (B_E)_n\to0.
\]
左端の単射は$\operatorname{Hom}_E(L_\bullet,h^nL_\bullet[-1])=0$、
右端の全射は$\operatorname{Ext}^1(P,s^{n-1}P)=0$による。
$n=0$では$B_0\simeq (B_E)_0$（directed categoryの定義および
次数とブロック内の直交性）。
自然変換$\epsilon$はhelixの対象上で$s$と可換である。
実際、$\epsilon_X$の自然性から
$\epsilon_X\circ s^{-1}\epsilon_X
 =\epsilon_X\circ\epsilon_{s^{-1}X}$となるが、
$X=s^aP$のとき後合成$\epsilon_X$は上と同じ負Ext消滅により
$\operatorname{Hom}(s^{-2}X,s^{-1}X)$上単射である。
よって$z_P=\epsilon_{sP}:P\to sP$は中心一次元となり、
左右の正則性と上の完全列から$B/(z_P)=B_E$である。

$(L_\bullet,h)$は$\coh E$のample pairである。
$h^{-n}L_i$の次数が$-3n+\deg L_i$であり、固定したsheafへの
Homによる生成性とExt消滅が従う。具体的には固定sheafを固定した
負のample線束の直和の商として表示する。十分高次数の任意の線束は
大域生成かつ$H^1=0$であり（次数二以上なら大域生成、正なら$H^1=0$）、
商の大域生成と$H^1$消滅を得る。この議論は特異な既約曲線でも成立する。
このためArtin--Zhang再構成\cite[Theorem~4.5]{AZ1994}で
$B_E$は右noetherian、$\qgr B_E\simeq\coh E$である。
この応用では$B_E$の次数零部分は有限次元代数でよい。
$B$も正則中心一次元による持上げで右noetherianとなる：
$z_P$-進随伴次数環は$B_E[Z]$であり、斉次右イデアルの初期項を
有限個持ち上げると、$z_P$-位数を増やす除法が固定全次数内で停止する。
斉次右イデアルの有限生成から通常の右イデアルの有限生成も従う。

$\qgr B$の任意の固定対象$X$は、AS正則性による有限生成射影加群の
有限分解から、helix対象の有限直和による有限分解を持つ。
従ってvery strongnessから
$\operatorname{Ext}^{>0}(s^{-n}P,X)=0$ ($n\gg0$)。
また十分負のhelix対象がheartを生成するので
$(P,s)$はample pairである。
ここでは標準の導来比較を用いており、dg商を無条件に通常の
アーベル圏の導来圏と同一視しているのではない。

任意の$M\in\coh E$について
\[
 \operatorname{Hom}(s^{-n}P,JM[k])
 =\operatorname{Hom}_E(h^{-n}L_\bullet,M[k])
\]
は$n\gg0$で$k\ne0$なら零である。
一方固定有界複体に対し上記のExt消滅を有限個のcohomology sheafに
適用すれば、左辺は
$\operatorname{Hom}(s^{-n}P,H^k(JM))$となる。
生成性により$JM$はheartに属し、$J$は完全である。
同じ等式と曲線上の生成性から$JM=0$なら$M=0$。
$JM$の切断加群は高次数で
$\bigoplus_n\operatorname{Hom}_E(L_\bullet,h^nM)$に一致し、
$B_E$を通じて作用する。実際$F\epsilon=0$は
$F\to FJF$がadjunctionにより分裂単射であることから従う。
従って$JM$は$\qgr(B/(z_P))\simeq\coh E$の像である。
この像の対象の部分対象・商対象も中心元で消される：
$\qgr$での中心元の作用が零なら切断加群の十分高次は
その中心元で消され、逆も成り立つ。
よってGabriel--Krull filtrationにおいてもこの閉部分圏は
次元一以下である。以上から主張を得る。
\end{proof}

\begin{lemma}[AZ2による特異点での局所自由性]
\label{lem:ou-az2-singular-free}
$E$を連結被約射影曲線、$h=\alpha^*(-)\otimes N$を幾何的自己同値とする。
$\coh E$の対象$M$がAZ2を満たす、すなわち任意の全射$V\to W$に対し
$n\gg0$で$\operatorname{Hom}(h^{-n}M,V)\to
\operatorname{Hom}(h^{-n}M,W)$が全射であるとする。
このとき$M$は全ての特異点で局所自由であり、
$M=V_0\oplus T$（$V_0$はベクトル束、$T$は滑らかな点に台を持つ
有限長層）と分解する。
\end{lemma}
\begin{proof}
特異点の集合は有限なので、任意の特異閉点$p$は$\alpha$の周期点である。
全射$\mathcal O_E\to k(p)$にAZ2を適用し、$n$を周期の十分大きい倍数に取る。
$H=(h^{-n}M)_p$、$A=\mathcal O_{E,p}$、$\mathfrak m$をその極大イデアルとする。
$\operatorname{Hom}_E(h^{-n}M,k(p))$は
$\operatorname{Hom}_A(H,A/\mathfrak m)$に等しいので、
大域的全射性から局所的にも
\[
 \operatorname{Hom}_A(H,A)\longrightarrow
 \operatorname{Hom}_A(H,A/\mathfrak m)
\]
が全射である。
$\mu=\dim_k H/\mathfrak m H$と置き、双対基底を$\mu$本の$A$値写像に
持ち上げると$H\to A^\mu$を得る。剰余体上では同型なので
Nakayamaの補題により全射であり、終域の自由性により分裂する。
核は剰余体上零となるので再びNakayamaの補題により零である。
従って$H$は自由である。$\alpha^n(p)=p$なので$M_p$も自由である。

$T$を$M$の最大零次元部分層とする。
被約曲線の局所自由層には零次元部分層がないので、
$T$の台は滑らかな点に限られる。
$V_0=M/T$は特異点で$M$と同じ、滑らかな点では
DVR上のtorsion-free加群なので局所自由である。
$\operatorname{Ext}^1_E(V_0,T)=H^1(E,V_0^\vee\otimes T)=0$
から分裂する。連結性により$V_0$のrankは一定である。
\end{proof}

\begin{proposition}[曲線上の反標準制限のrank認識]
\label{prop:ou-reduced-curve-rank}
$E$を連結被約射影算術種数一曲線、$h=\alpha^*(-)\otimes N$、
$\deg N=3$とする。$\mathcal C$の反標準移動$s$と内在的三次標識
$\mathcal O$に対し、導来随伴$F\dashv J$、
自然な三角$s^{-1}\xrightarrow{\epsilon}\mathrm{id}\to JF$、
$Fs\simeq hF$があり、$J$はheart上完全かつconservativeで、
その連接層の像がGK次元一以下であるとする。
このとき$M=F\mathcal O$は線束である。
\end{proposition}
\begin{proof}
$\epsilon_{\mathcal O}$はheart内の射であり、そのconeが$JF\mathcal O$。
$J$の完全性とconservativityから$F\mathcal O$のcohomologyは
次数$-1,0$に限られ、$JH^{-1}(F\mathcal O)$は
$s^{-1}\mathcal O$の次元一以下の部分対象となる。
完全自己同値$s$はGabriel--Krull次元を保つので、
$s^{-1}\mathcal O$のtorsion-free性からこれは零となる。従って$M=F\mathcal O\in\coh E$で
\begin{equation}\label{eq:ou-O-divisor}
 0\to s^{-1}\mathcal O\to\mathcal O\to JM\to0.
\end{equation}
任意の全射$V\to W$ in $\coh E$は$JV\to JW$という全射を与え、
AZ2と随伴により$\operatorname{Hom}_E(h^{-n}M,V)\to
\operatorname{Hom}_E(h^{-n}M,W)$が$n\gg0$で全射となる。
補題\ref{lem:ou-az2-singular-free}により$M=V_0\oplus T$、
$V_0$はrank $r$のベクトル束、$T$の長さを$\ell$とする。
$r>0$である：任意の非零線束$V$について$JV\ne0$であり、
AZ1による生成性と随伴から、ある正の$n$について
$\operatorname{Hom}_E(h^{-n}M,V)\ne0$となるが、
零次元層から線束への射は零である。

AZ2を(\ref{eq:ou-O-divisor})の全射に適用すると$n\gg0$で
\[
 \dim\operatorname{Hom}_E(M,h^nM)
 =\dim\operatorname{Hom}_{\mathcal C}(s^{-n}\mathcal O,JM)
 =\dim R_n-\dim R_{n-1}=3n .
\]
$E$の全次数で$\deg(h^nV_0)=\deg V_0+3nr$だから、
曲線のRiemann--Rochにより
$\chi_E(V_0,h^nV_0)=3nr^2+r^2\chi(\mathcal O_E)$。
特に算術種数一なら
\[
 3n\geq 3nr^2+r\ell.
\]
従って$r=1,\ell=0$である。
\end{proof}

\begin{proof}[定理\ref{thm:ou-integral-elliptic-recognition}の証明]
補題\ref{lem:ou-elliptic-shift}、\ref{lem:ou-elliptic-heart}により
命題\ref{prop:ou-reduced-curve-rank}の全仮定を満たす。
従って$M=F\mathcal O$は線束で、完全列(\ref{eq:ou-O-divisor})が成り立つ。
既約種数一曲線では正次数の線束の$H^1$は零だから、
任意の$n\geq1$で$\dim\operatorname{Hom}_E(M,h^nM)=3n$。
(\ref{eq:ou-O-divisor})にHomを施した左完全列と次元比較から、
制限写像は全ての正次数で全射となる。次数零も恒等元により全射。
$\epsilon$は全ての$s^a\mathcal O$上単射なので、自然性から
上述の$P$の場合と同様に$s$と可換である。
従って$z=\epsilon_{s\mathcal O}$は正則中心一次元で
\begin{equation}\label{eq:ou-R-mod-z}
 R/zR\simeq \bigoplus_{n\geq0}\operatorname{Hom}_E(M,h^nM)
 \simeq B(E,M^\vee\otimes hM,\alpha)
\end{equation}
となる。右辺はorbit積の規約による自己同型（反対の規約では逆向き）
のtwisted homogeneous coordinate ringである。

最後に環論的認識を記す。これは\cite[Section~12]{VdBBlowup2001}の
中心拡大の議論を用いるが、同論文の爆発表示を仮定しない。
$\bar R=R/zR$の偏極$\mathcal L=M^\vee\otimes hM$は次数三であり、
既約Gorenstein種数一曲線上でvery ampleである。
実際、長さ二の任意の閉部分スキーム$Z$に対して
$H^1(\mathcal L\otimes I_Z)=0$：Serre双対性により双対空間は
$\operatorname{Hom}(\mathcal L\otimes I_Z,\mathcal O_E)$だが、
正次数のrank一torsion-free層から$\mathcal O_E$への非零射は
次数比較で不可能である。従って$|\mathcal L|$は$E$を平面三次曲線に埋め込む。
さらに$\alpha^*$は$\Pic^0(E)$上恒等なので
\[
 \mathcal L\otimes(\alpha^{2})^*\mathcal L
      \simeq(\alpha^*\mathcal L)^{\otimes2}.
\]
これはATV\cite[(4.8), Theorems~6.6, 6.8]{ATV1990}のregular triple条件。
ATVのelliptic tripleの定義は$E$の滑らかさを要求しない。
従って$\bar G=Q(\bar R)$は三生成元の二次AS三次元正則環であり、
$\bar R$はその正則正規三次商である。
偏極を保つ場合も同じ定理に含まれる。正則な$z$と$\bar R$の一次生成性により
$R$は一次生成。$G=Q(R)$とし、$W=R_1$、
$S_2=\ker(W^{\otimes2}\to R_2)$とすれば$\dim W=4,\dim S_2=6$。
$z$は$G$でも中心元で$G/zG=\bar G$である。
$\dim G_3\leq\dim G_2+\dim\bar G_3=20$。
$\bar G$の三次関係の持上げは$R$への核に入り、
$\bar G$上で非零なので$G$でも非零である。
（具体的に$\bar R$で消える三次式を持ち上げ、その$R$上の像を
$zR_2$で補正する。）$\dim R_3=19$より$\dim G_3=20$、
よって$z:G_2\to G_3$も単射である。

中心拡大判定\cite[Corollary~2.7]{LBSVdB1996}の条件を確認する。
$z:G_1\to G_2$の単射は$R$での単射から従う。
$\bar W=W/kz$、$\bar S_2$を$\bar G$の二次関係空間とする。
$\bar W\otimes\bar S_2\cap\bar S_2\otimes\bar W$の任意の元を
$u\in W\otimes S_2$、$v\in S_2\otimes W$にそれぞれ持ち上げる。
差$u-v$は$z$を含むテンソルの和であり、隣接交換関係
$C=\langle z\otimes w-w\otimes z:w\in W\rangle\subset S_2$
を法として$z\otimes q$ ($q\in W^{\otimes2}$)と書ける。
$G_3$で$zq=0$かつ$z:G_2\to G_3$が単射なので$q\in S_2$。
従って$u-v$は$z\otimes S_2+W\otimes C+C\otimes W$に属し、
各項は射影で消える。これらを両持上げから補正すれば
同一の元$W\otimes S_2\cap S_2\otimes W$を得る。
従って線形syzygyの持上げも成立する。
中心拡大判定から$G$はAS四次元正則、Hilbert級数$(1-t)^{-4}$。
$G/zG$のHilbert級数比較から$z$は$G$上でも正則であり、
$\bar G$の左右noetherian性を持ち上げられる。
また$\bar G$はKoszulで、正則中心一次元で割った環への
change-of-rings Tor spectral sequenceは対角次数に集中する
（$\operatorname{Tor}^G_q(\bar G,k)$は$q=0,1$のみ、
内次数$q$）。従って$G$もKoszulである。
命題\ref{prop:quadratic-envelope-recognition}から$R=G/(g)$。
AZ再構成によりこれは与えられた標識と反標準自己同値に適合する。
\end{proof}

\begin{proposition}[被約可約三次曲線への条件付き拡張]
\label{prop:ou-reduced-cubic-geometric}
$E$を連結被約平面三次曲線とし、$F,J,h$が
命題\ref{prop:ou-reduced-curve-rank}の仮定を満たすとする。
さらに$h$が$\Perf E$上で次数$0,1,2$の各三本の線束$L_i$に沿う
spherical twistの積の逆を$[2]$したものとする。
このとき与えられた内在的標識の反標準環$R$は
標準AS四次元正則環の正則正規三次商である。
\end{proposition}
\begin{proof}
命題\ref{prop:ou-reduced-curve-rank}から$M=F\mathcal O$は線束である。
まず自己同型$\alpha$は全ての既約成分を個別に保つ。
実際、任意の全次数零の線束$Q$に対し、
$[Q]-[\mathcal O_E]\in K_0(\Perf E)$は各twistで不変である。
曲線Riemann--Rochから
$\chi(L_i,Q)-\chi(L_i,\mathcal O_E)=\deg Q=0$だからである。
幾何的な$h$のもとで行列式を取れば$\alpha^*Q\simeq Q$。
各成分の滑らかな点の差で任意の全和零の多重次数を実現できるので、
成分の置換はこの多重次数格子に恒等として作用する。
従って置換自体が恒等である。
この議論は$\Pic^0(E)$のみの作用より強い。

さらに$(M,h)$はAZ1も満たす。任意の$V\in\coh E$に対し
AZ1で全射$\bigoplus_j s^{-n_j}\mathcal O\to JV$を取る。
随伴で対応する$\bigoplus_j h^{-n_j}M\to V$が全射でなければ、
$J$の完全性により元の全射は$J(\operatorname{im})$を経由し、
$J(\operatorname{coker})\ne0$と矛盾する。
成分$E_i$への制限で$\deg N|_{E_i}\leq0$なら、
$h^{-n}M|_{E_i}$の次数は$n>0$で下がらない。
従って十分負次数の線束を$E_i$から押し出した対象への射は全て零であり、
AZ1に反する。よって全成分で$\deg N|_{E_i}>0$である。
$\mathcal L=M^\vee\otimes hM$も同じ正の多重次数を持ち、全次数三である。

このような$\mathcal L$はvery ampleである。証明のため
任意の長さ二の零次元閉部分スキーム$Z\subset E$を取る。
Serre双対性で$H^1(\mathcal L I_Z)^*$は
$\operatorname{Hom}(\mathcal L I_Z,\mathcal O_E)$に等しい。
非零射の像が一般点で非零となる成分の和を$E'\subset E$とする。
$E'=E$なら射は単射であり、源の次数$3-2=1$は終域の次数零より
大きいので矛盾する。
$E'$が真の部分曲線なら、終域のうち$E'$に台を持つ最大の部分層は
$\mathcal O_{E'}(-E'\cap\overline{E\setminus E'})$である。
これは平面内で補曲線の方程式を掛けることで得る。
交叉長はBezoutの定理により二である（次数$1$と$2$）。
射は源を$E'$に制限しその零次元torsionを除いた
$I_Z\mathcal O_{E'}\otimes\mathcal L|_{E'}$を経由し、
この層から上の終域への単射を与える。
源の次数は$\deg(\mathcal L|_{E'})-2\geq-1$、
終域の次数は$-2$であり、再び矛盾する。
従って$H^1(\mathcal L I_Z)=0$で、長さ二の部分スキームを分離する。
同じ支持部分曲線の議論で、全成分上正次数の任意の線束の$H^1$は零となる。
よって$h^n$のorbitに対しても
$\dim\operatorname{Hom}(M,h^nM)=3n$ ($n\ge1$) を得る。

完全列(\ref{eq:ou-O-divisor})と次元比較から、
定理\ref{thm:ou-integral-elliptic-recognition}の証明と同じく
正則中心一次元$z$と$R/zR=B(E,\mathcal L,\alpha)$を得る。
$\alpha^*$は全次数零線束を固定するので
$\mathcal L\otimes(\alpha^2)^*\mathcal L\simeq
(\alpha^*\mathcal L)^{\otimes2}$。
ATVのregular triple条件と同定したvery ample性から
$Q(R/zR)$は三次元二次AS正則環である。
以後は定理\ref{thm:ou-integral-elliptic-recognition}の環論的部分と
全く同じ中心拡大判定により結論を得る。
この判定で$R/zR$の整域性は用いず、$R$上の$z$の正則性と
三次関係の持上げを用いるため、可約曲線でも適用できる。
\end{proof}

\begin{remark}[可約曲線での実際の残存課題]\label{rem:ou-reduced-curve-gap}
命題\ref{prop:ou-reduced-cubic-geometric}は幾何的$h$と
曲線heartの認識を仮定している。既約の場合の
\cite[Corollary~5.8]{BurbanKreu2006}は可約曲線にそのまま使えない。
ただし、被約平面三次曲線上の九線束が同一ブロック内完全直交、
異なるブロック間前向きhigher Ext消滅のfoundationを与え、
標準noetherian OU helix/導来比較が成立する場合に、
$h=\alpha^*(-)\otimes N$で$N$がampleであることを別途証明できれば、
補題\ref{lem:ou-elliptic-heart}の証明はそのまま拡張できる。
実際、$N$の全成分で正の次数から負のhelix像による
大域生成と高次消滅を得る。十分高い多重次数の線束については
上の部分曲線への制限とSerre双対性による議論で
$H^1$消滅と長さ一の部分スキーム上の全射性が従う。
任意の連接層については固定した負のample線束による全射から従う。
degree zeroでの$B_0\simeq(B_E)_0$はfoundationの仮定による。
従って$J$の完全性・conservativity・像の次元一以下を確認できる。
次の定理\ref{thm:ou-reduced-md-criterion}は、このような幾何的かつampleな$h$の存在を九線束の多重次数により必要十分に判定する。
その条件外についても節点型では定理\ref{thm:ou-nodal-complete-recognition}により
標準OU helix比較を持つ九線束表示の全てを認識した。
加法型についても定理\ref{thm:ou-additive-partial-recognition}により
III型は全候補を認識した。さらに系\ref{cor:ou-iv-six-residual}により
IV型の最後の6候補を一つの明示的関係代数族に帰着し、
定理\ref{thm:ou-iv-periodic-exclusion}によって内在的三次条件から排除した。
従って定理\ref{thm:ou-reduced-complete-recognition}により全被約平面三次曲線の
九線束表示について認識が完了する。
非被約平面三次曲線についても
定理\ref{thm:ou-nonreduced-plane-recognition}は既約尖点曲線への帰着を与える。
非被約曲線では特異点集合の有限性を使った
補題\ref{lem:ou-az2-singular-free}を適用せず、
合成の生成性と縮約を用いる。
従って定理\ref{thm:ou-all-plane-cubic-recognition}により
平面三次曲線による九線束表示全体を認識する。
一般のOU型に対する曲線表示の存在は未解決であり、
定理\ref{thm:ou-smooth-determinant-reconstruction}がそのための
具体的な十分条件を与える。
\end{remark}

\begin{remark}\label{rem:ou-smooth-elliptic-scope}
この定理は指定したOU helix比較をもつ既約種数一曲線型の
内在的候補に対する肯定的認識であって、曲線データだけから
内在的条件を導いたものではない。爆発中心の一般位置、
関係係数の全非消失、平行移動の非torsion性はいずれも
認識定理の追加条件ではない。
特に既知の非可換爆発比較の適用外でも、上の仮定を満たす限り
標準超曲面表示を得る。可約被約型の肯定的範囲は系\ref{cor:ou-reduced-md-recognition}でさらに広がる。
一般のOU型がこれらの曲線表示を持つか、平面三次曲線表示を持たない型で
同じ認識を実現できるかは残る問題である。
spherical functorと中心拡大判定自体は既知であり、
ここでの新しい主張は内在的ample条件による制限線束の認識と
それらの接続である。新規性の網羅的確認は未了。
\end{remark}

\subsection{射影加群の切断によるOU比較と既約三次曲線型}
\label{subsec:ou-perfect-tail-comparison}

OU\cite[\S2]{OU}の $\operatorname{mod} A$ はnoetherianな右加群の圏である。
定義1.1(2)は$qgr$表示を要求し、定義1.1(3)は導来圏内の標識を指定する。
ここでは前者から出発する。OUの(3.9)--(3.10)に対応する計算をまず
射影加群の三角圏で行うことで、右noetherian性を曲線から証明できる。

\begin{proposition}[OU型の有限次元加群による商]
\label{prop:ou-perfect-tail-comparison}
$A$をOUの意味でAS正則、次元三、parameter $\ell$とし、
単純右加群がOU定義5.2の有限分解(5.4)を持つとする。
$P_i=e_i A$, $S_i$をその単純商とし、
\[
 \mathcal D=\big(\Perf A/\operatorname{thick}(S_i\mid i\in\mathbb Z)\big)^{\natural}
\]
と置く（$\natural$は冪等完備化）。$P_i$の像を$E_i$とする。
任意の長さ$\ell$の窓は$\mathcal D$のfull strong exceptional collectionであり、
$\operatorname{Hom}(E_i,E_j)=A_{ji}$、$E_{i+\ell}\simeq S_{\mathcal D}^{-1}E_i[2]$
である。全列はvery strongである。窓の自己準同型代数を$\Lambda$とすると
$\mathcal D\simeq\Perf\Lambda$である。
各$P_i$が右noetherianならさらに標準の同値
$\mathcal D\simeq D^b(\qgr A)$が成立し、$E_i$は$\pi P_i$に移る。
\end{proposition}
\begin{proof}
右加群の添字が作用で減少する規約を取る。
$P_{i,\le m}$を添字$m$以下への切断、$T_m=P_i/P_{i,\le m}$とする。
$T_m$は有限長で、各$S_a$は(5.4)によって完全なので切断も完全である。
局所化の射は
\[
 \operatorname{Hom}_{\mathcal D}(E_i,E_j[d])
 =\varinjlim_{m\to-\infty}\operatorname{Ext}^d_A(P_{i,\le m},P_j)
\]
で計算できる。実際、有限個の$S_a$から作った任意の有界対象$K$と
全てのシフトに対し、十分小さい$m$で
$\operatorname{Hom}(P_{i,\le m},K[d])=0$となる。
切断は添字$m$以下の射影のみで分解できるからである
（各段で台の中に斉次生成元を取る）。従って錐がこの部分圏に入る
任意の分母$X\to P_i$を十分深い切断が経由する。同じ消滅を錐の
シフトに使えば二つの持上げの差もさらに切断して消える。
これはVerdier商のroofの余終性を与える。

OU (2.1)と$T_m$の組成列から、
$\operatorname{Ext}^d_A(T_m,P_j)$は$d=3$以外で零、次数三の次元は
$(T_m)_{j+\ell}$の次元に等しい。この次数三の群は$m<j+\ell$
で安定する（隣接する切断の差の組成因子に再び(2.1)を用いる）。
切断の短完全列から次数零のHomは射影間の射そのものであり、
合成も保たれる。残りは次数二のみで、その次元は$\dim A_{i,j+\ell}$。
従って各窓の例外性・強性と全列のvery strong性が従う。
(5.4)の両端は$P_{i-\ell},P_i$、中間項の添字はその間にある。
商で$S_i=0$なので、窓の生成部分圏は左右に一項ずつ延長できる。
よって窓はfullであり、dg商のtilting対象が$\mathcal D\simeq\Perf\Lambda$
を与える。$\Lambda$は有限次元有向代数なので有限大域次元を持つ。

窓$E_{i+1},\ldots,E_{i+\ell}$の直和を$T$とする。
上のHom計算による$\operatorname{RHom}(T,E_i[2])$と
Serre双対性による$\operatorname{RHom}(T,S_{\mathcal D}E_{i+\ell})$
はともに末頂点の一次元単純加群を次数零に置いたものなので周期公式を得る。

最後に右noetherian性を仮定すると、有限生成加群は有限生成射影による
有限分解を持ち、$\Perf A=D^b(\operatorname{mod} A)$である。
この有界導来圏で有限次元cohomologyを持つ部分圏は上記の
$\operatorname{thick}(S_i)$である。これによる商を$D^b(\qgr A)$
と比較するには、アーベル商の射を十分深い切断間の射で表す。
実際、有限生成torsion加群は有限次元であり、Serre商で用いる
始域の有限余次元部分加群・終域の有限次元核は深い切断で消去できる。
有界複体、chain map、homotopyの各有限図式は共通の切断を取れば
持ち上がる。積が零、または二射が等しいという有限個の条件も同様である。
$qgr$で擬同型となる写像の錐は有限次元cohomologyを持つ。
従って有界複体のroofとその等価関係も持ち上がり、所望の同値を得る。
有界導来圏は冪等完備である。
\end{proof}

\begin{corollary}[既約曲線表示からの標準比較]
\label{cor:ou-integral-automatic-comparison}
$A$をAOU三ブロック箙型のOU代数とする。九頂点の窓の有向代数が、
定理\ref{thm:ou-integral-elliptic-recognition}の既約曲線上の九線束の
有向代数と頂点付きで同型なら、各$P_i$は右noetherianであり、
標準OU比較が成立する。従って$\qgr A$に内在的三次標識があれば、
その反標準環は標準AS四次元正則環の正則正規三次商である。
\end{corollary}
\begin{proof}
まず命題\ref{prop:ou-perfect-tail-comparison}の$\mathcal D$の中で
$P=\bigoplus_{i=1}^9 E_i$, $s=S_{\mathcal D}^{-1}[2]$と置く。
$F,J,h$の構成と補題\ref{lem:ou-elliptic-shift}は
$\Perf\Lambda$と曲線だけを用いる。補題\ref{lem:ou-elliptic-heart}
の証明のうち、軌道環$B$の右noetherian性までの議論をここで適用できる。
実際、$n\ge1$の短完全列
$0\to B_{n-1}\to B_n\to\operatorname{Hom}_E(FP,h^nFP)\to0$
の左端は曲線上の負Ext消滅、右端は$\mathcal D$の窓の強性と
very strong性による。$n=0$は有向代数の同定による。
同補題の自然性と単射性の計算から核は正則中心一次元で生成される。
曲線上の$h$は次数三の幾何的移動であるから商環はAZの定理により
右noetherianであり、中心元による持上げで$B$も右noetherianとなる。
この段階で$\mathcal D$上のheartは必要ない。

九つの合同類ごとに周期同型を反復して選ぶと、$B$を九頂点と周期で
展開した代数の各成分と合成は、同命題の次数零Homの同定により
$A$そのものである。逆向きの周期には正次数付き性により射はない。
各$P_i$の部分加群はこの展開を戻すと、$B$の頂点射影の次数シフトの
斉次部分加群に対応する。従って全$P_i$が右noetherianになる。
同命題の最後の主張をここで初めて用い、既約曲線認識定理を適用する。
\end{proof}

\begin{corollary}[正の多重次数を持つ曲線表示での比較]
\label{cor:ou-positive-automatic-comparison}
同じOU三ブロック窓が注意\ref{rem:ou-reduced-curve-gap}の被約平面
三次曲線上の九線束で表示され、曲線上の移動が
$h=\alpha^*(-)\otimes N$、全成分で$\deg N>0$であれば、
標準OU比較は成立する。特に定理\ref{thm:ou-reduced-md-criterion}の
$\delta_{\rm md}=0$かつ全$w_i>0$の範囲では、内在的標識の反標準環が
標準三次商になる。
\end{corollary}
\begin{proof}
前系の証明で既約性を用いたのは$h$の幾何性・豊富性を導く部分である。
ここではこれらが与えられる。次数零での窓のHomの同定には、前向き
高次Extの消滅とGorenstein曲線のSerre双対性から得る逆向きHomの消滅を使う。
正の多重次数に対する注意\ref{rem:ou-reduced-curve-gap}の生成性・消滅により
$(FP,h)$はampleであり、曲線の軌道環は右noetherianである。
従って$\mathcal D$での中心拡大、展開、右noetherian性を経て
命題\ref{prop:ou-perfect-tail-comparison}が適用される。
比較を得た後は命題\ref{prop:ou-reduced-cubic-geometric}を用いる。
\end{proof}

\subsection{既約種数一曲線による任意のOU三次箙の認識}
\begin{theorem}[箙を指定しない既約種数一曲線窓]
\label{thm:ou-any-quiver-elliptic-window}
$A$を任意のOU三次型のAS正則$\mathbb Z$代数とする。
命題\ref{prop:ou-perfect-tail-comparison}の九頂点の射影窓の有向dg圏が、
滑らかな楕円曲線、既約節点三次曲線または尖点三次曲線$E$上の九つの完全な球面的対象
$V_1,\ldots,V_9\in D^b(\coh E)$の有向dg圏と、頂点を保って準同値とする。
有向dg圏では前向きの全導来Homを取り、逆向きの射を零、自己準同型を
恒等射のスカラー倍に制限する。このとき標準射影は右noetherianで、
$\qgr A\in\mathfrak I_3$なら任意の内在的標識について
標準四次元AS正則環の正規三次商の表示を得る。
曲線上の九対象の階数・次数・シフト、箙の三ブロック性、
標準OU比較を仮定しない。
\end{theorem}
\begin{proof}
$P_i$を完全商$\mathcal D$の窓、$s=S_{\mathcal D}^{-1}[2]$とする。
有向dg圏による球面的構成\cite[(3)--(5)]{OUSpherical2020}は
$F(P_i)=V_i$となる関手$F$、右随伴$J$と三角
$s^{-1}\to\mathrm{id}\to JF$を与える。
$Fs=hF$、$h=(T_{V_1}\cdots T_{V_9})^{-1}[2]$は
式\eqref{eq:ou-restriction-commute}と同じ随伴公式による。特異曲線でもこの構成は\(\Perf E\)のSerre関手$[1]$を使い、$J$は完全な$V_i$からの導来Homにより全ての有界連接複体上に延長する。

OU定義1.1(2)の同じ箙を与える滑らかな三次del Pezzo曲面を$X$、
その窓を$U_i$とする。共通のシフトを除き$U_i$は全て正階数の
ベクトル束である。実際、任意の$i,j$について
$R\mathcal H om_X(U_i,U_j)$を$-nK_X$で捻った超cohomologyは
$n\gg0$で次数零に集中する（very strong性）。
Serre消滅と大域生成からこの層複体自体のcohomologyが次数零に集中する。
特に各点の局所環で$U_i$の最小有界自由複体を取り、
両端の次数を$a<b$とすれば、その自己Hom複体の最上次数$b-a$への
微分の像は極大イデアル倍に入るため、そのcohomologyは非零となり矛盾する。
従って局所的に一つの次数だけで自由である。台は開かつ閉で、連結性から
台は全体、シフトも一定である。異なる$i,j$にも同じ消滅を使えばシフトは共通である。
$r_i=\mathrm{rk}(U_i)>0$、$d_i=-K_X\cdot c_1(U_i)$をこの共通シフトを除いて取る。

完全商と古典窓は同じEuler行列を持つ。実際、OU (5.4)の固定した
射影分解は各端点でのHilbert次元を対角の$1$とそれより外側の$0$から
添字差の順に一意に定める。窓上では命題\ref{prop:ou-perfect-tail-comparison}
によりこれが前向きHom次元となる。窓がfull exceptionalなので
この基底同士の同型はEuler形式を保ち、Serre作用も保つ。
$v_i=[V_i]_{\rm num}\in\mathbb Z^2$を階数・次数の順に取ると、
上の三角のEuler特性を古典窓で計算して
\[
 \det(v_i,h^n v_j)
  =\chi(P_i,s^nP_j)-\chi(P_i,s^{n-1}P_j)
  =r_i d_j-d_i r_j+3n r_i r_j .
\]
最後の等式は$X$の滑らかな反標準切断へ制限する完全列と
曲線Riemann--Rochによる。$(r_i,d_i)$は$\mathbb Z^2$を生成する。
実際、full exceptional基底の整数結合として$[\mathcal O_X]$と
一本の$(-1)$曲線$D$の$[\mathcal O_D]$を取り、
階数と反標準次数を取れば標準基底になる。同じ整数結合を$v_i$に施すと、
$n=0$の等式により行列式一の基底を得る。各$v_i$のこの基底での座標も
二つの行列式により$(r_i,d_i)$に一致する。
これらの既約曲線の導来自己同値は$\mathrm{SL}_2(\mathbb Z)$上に全射する
\cite[Corollary~5.8]{BurbanKreu2006}。そこで$F$に一つの自己同値を後合成して
$v_i=(r_i,d_i)$とする。$n=1$の等式から$h$の数値作用は
$(r,d)\mapsto(r,d+3r)$である。

同じ自己同値の記述により$h=\alpha^*(-)\otimes N[2k]$、
$\deg N=3$である。$k\ne0$なら大きな$n$では
$\operatorname{Hom}^p(V_i,h^nV_j)=0$ ($p=-1,0$)である。
実際、線束・自己同型を作用させた固定有界複体の振幅は一様有界で、
各$V_i$は完全であり、その双対とのテンソル積の超cohomologyも一様有界だからである。三角から
$\operatorname{Hom}(P_i,s^{n-1}P_j)\simeq\operatorname{Hom}(P_i,s^nP_j)$
となり、very strong性と上の正のEuler差分に反する。従って$k=0$。
Burban--Kreu{\ss}ler\cite[Proposition~4.13]{BurbanKreu2006}の
既約種数一曲線上の球面的対象の分類と正階数から、
$V_i=W_i[k_i]$、$W_i$は安定ベクトル束、$k_i$は偶数である。
三角とvery strong性から$n\gg0$で
$\operatorname{Hom}^p(V_i,h^n V_j)=0$ ($p\notin\{-1,0\}$)。
Riemann--Rochにより$\operatorname{Hom}(W_i,h^nW_j)\ne0$なので
$k_i-k_j\in\{-1,0\}$、従って全シフトは等しい。共通シフトを除く。

$W=\bigoplus_i W_i$とする。$(W,h)$は$\coh E$のample pairである。
非常に豊富な線束$H$を固定する。$h^{-n}W_i$と
$h^{-n}W_i\otimes H$は傾き$-3n+O(1)$の安定束である。
固定層$M$を固定線束の有限直和の商に表せば、どちらの束への$M$からの射も
十分大きな$n$で消える。Gorenstein曲線のSerre双対性により
$(h^{-n}W_i)^\vee\otimes M$とその$H^{-1}$によるtwistはともに
$H^1=0$である。従って前者はCastelnuovo--Mumford正則性により大域生成される。評価写像と
正階数の局所自由性から$h^{-n}W_i$の有限直和が$M$に全射する。

$P=\bigoplus_i P_i$の周期環を$B$とする。
補題\ref{lem:ou-elliptic-heart}の短完全列は再び成立する。
正次数での単射には曲線上の負Ext消滅、全射には窓のvery strong性を用いる。
次数零では$\operatorname{Hom}^p(P,s^{-1}P)=0$ ($p=0,1$)が
Serre双対と窓の強性から従う。ゆえに$B$は$(W,h)$の切断環の
正則中心一次拡大である。切断環の右noetherian性と中心元による持上げ、
九頂点への展開から、系\ref{cor:ou-integral-automatic-comparison}の
証明により全標準射影は右noetherianとなり、ここで標準比較を得る。
また同補題の高twistの随伴Hom判定により$J$はheart上で完全かつ
conservativeであり、その像は中心元で消されて次元一以下となる。
命題\ref{prop:ou-reduced-curve-rank}から内在的標識の$F$像は線束である。
$h$に現れる各球面的twistはEuler形式の根基に恒等として作用する。
$[Q]-[\mathcal O_E]$ ($Q\in\Pic^0 E$)に行列式を適用すれば
$\alpha^*|_{\Pic^0 E}=\mathrm{id}$である。
ここから式\eqref{eq:ou-R-mod-z}およびそれに続く中心拡大の証明を
そのまま適用して所望の表示を得る。
\end{proof}

特にOU\cite[Definition~6.3, Theorem~6.4]{OU}の有理対応が
滑らかな楕円曲線上の球面的列から窓を復元する一般ポテンシャルの範囲には、
この定理が各三次箙で適用できる。有理対応の不定点や
滑らかな楕円曲線による窓をまだ持たない特殊化には別の議論が必要である。
